\title{DeepSeekMath: Pushing the Limits of Mathematical Reasoning in Open Language Models}

\begin{abstract}
Mathematical reasoning remains a formidable obstacle for language models owing to its inherent complexity and highly structured characteristics. This study presents DeepSeekMath~7B, an extension of DeepSeek-Coder-Base-v1.5 7B, further pre-trained on 120B math-specific tokens extracted from Common Crawl, supplemented by both natural language and code corpora. Without recourse to external toolkits or ensemble strategies, DeepSeekMath~7B attains a notable 51.7\% on the MATH competition benchmark, coming close to the results of Gemini-Ultra and GPT-4. When utilizing self-consistency across 64 generations, DeepSeekMath~7B reaches a score of 60.9\% on the MATH dataset. The model’s advancement in mathematical reasoning is credited to two principal innovations: first, a sophisticated data selection system leverages extensive, publicly available web resources; second, the application of Group Relative Policy Optimization (GRPO)—an adaptation of Proximal Policy Optimization (PPO)—which boosts the model's reasoning proficiency while also enhancing PPO's memory efficiency.
\end{abstract}

\section{Introduction}

Large language models (LLMs) have transformed the landscape of mathematical reasoning within artificial intelligence, driving progress in quantitative benchmarks~\citep{MATH} as well as geometry reasoning tasks~\citep{trinh2024solving}. Additionally, these systems have demonstrated effectiveness in supporting human users confronted with intricate mathematical challenges~\citep{tao}. Nonetheless, leading proprietary models like GPT-4~\citep{gpt4} and Gemini-Ultra~\citep{gemini} remain inaccessible to the public, and extant open-source alternatives exhibit a substantial performance gap.

This paper presents DeepSeekMath, a specialized language model in the mathematics domain that delivers mathematical reasoning and problem-solving performance superior to that of existing open-source models, nearing that of GPT-4 across established academic evaluations. The development of DeepSeekMath centers on constructing the DeepSeekMath Corpus, a substantial and high-quality pre-training dataset composed of 120 billion math-focused tokens. To compile this dataset, content is sourced from Common Crawl (CC) and identified with a fastText-based classifier \citep{joulin2016fasttext}. The initial version of the classifier is built using positive samples drawn from OpenWebMath \citep{openwebmath} and a wide array of unrelated web documents as negative samples. This classifier is subsequently employed to extract more positive data from the CC, which are then subjected to meticulous human annotation for further quality control. The annotated data are integrated to retrain the classifier, yielding improved extraction accuracy. Evaluation outcomes demonstrate the high fidelity of the constructed corpus; for instance, DeepSeekMath-Base 7B attains a 64.2\% accuracy on GSM8K \citep{gsm8k} and 36.2\% on the MATH competition benchmark \citep{MATH}, surpassing the scores of Minerva 540B \citep{minerva}. The multilingual composition of the DeepSeekMath Corpus additionally contributes to notable gains on Chinese mathematical assessments \citep{wei2023cmath,agieval}. We posit that our methodologies in mathematical dataset construction can serve as a reference point for the field, acknowledging considerable opportunities for enhancement in future research.

DeepSeekMath-Base is constructed by initializing with DeepSeek-Coder-Base-v1.5 7B \citep{deepseek-coder}, as empirical evidence suggests that commencing from a code-pretrained foundation yields superior results compared to leveraging a generic large language model. In addition, our findings reveal that mathematical pre-training not only enhances mathematical competencies but also leads to notable gains on benchmarks such as MMLU \citep{mmlu} and BBH \citep{bbh}, thereby augmenting general reasoning performance alongside domain-specific abilities.

Following the pre-training phase, DeepSeekMath-Base undergoes further refinement through mathematical instruction tuning, incorporating data reflecting chain-of-thought \citep{cot}, program-of-thought \citep{pot,pal}, and tool-integrated reasoning \citep{tora} approaches. The instruction-tuned model, DeepSeekMath-Instruct 7B, outperforms other 7B models and achieves performance on par with open-source models as large as 70B in parameter count.

In this work, we also present the Group Relative Policy Optimization (GRPO), a novel reinforcement learning (RL) method adapted from Proximal Policy Optimization (PPO) \citep{schulman2017proximal}. GRPO dispenses with the critic component, estimating the baseline by aggregating group scores, thus offering substantial computational efficiency. Utilizing only a selected portion of English instruction tuning data, GRPO achieves significant advancements over the already strong DeepSeekMath-Instruct, as evidenced by improvements in both in-domain tasks (GSM8K: 82.9\% $\rightarrow$ 88.2\%, MATH: 46.8\% $\rightarrow$ 51.7\%) and out-of-domain evaluations (such as CMATH: 84.6\% $\rightarrow$ 88.8\%) during RL fine-tuning. Additionally, we propose a comprehensive framework to contextualize various algorithms, including Rejection Sampling Fine-Tuning (RFT) \citep{yuan2023scaling}, Direct Preference Optimization (DPO) \citep{dpo}, PPO, and GRPO, and elucidate their characterization as either direct or streamlined RL strategies. Our rigorous experimental analysis—spanning aspects such as online versus offline learning, outcome versus process-based supervision, and the effects of single-turn versus iterative RL—allows for a deeper examination of this conceptual space. Finally, we analyze the mechanisms underlying the performance gains realized by RL on instruction-tuned models, and outline promising future directions for refining RL methods under this unified perspective.

\subsection{Contributions}
This work advances the field through large-scale mathematical pre-training and a thorough examination of reinforcement learning approaches.

\textbf{Math Pre-Training at Scale}

\begin{itemize}[topsep=0pt]
    \item We demonstrate that the widely available Common Crawl dataset holds substantial, untapped resources for mathematical model development. By crafting an advanced data curation pipeline, we compile the DeepSeekMath Corpus—a meticulously filtered, high-fidelity collection of 120B tokens sourced from mathematics-relevant webpages. This dataset greatly surpasses previous math corpora, being nearly 7 times larger than that used by Minerva \citep{minerva} and 9 times the scale of OpenWebMath \citep{openwebmath}.
    
    \item The DeepSeekMath-Base 7B, our pre-trained base model, achieves performance metrics on par with Minerva 540B \citep{minerva}, suggesting that sheer parameter count does not solely determine success in mathematical reasoning. High-caliber, focused pre-training on a more compact model can lead to similarly robust results.

    \item Insights from our training experiments reveal that conducting code pre-training before math-specific pre-training enhances the model’s problem-solving skills for mathematics, both with and without external computational tools. This provides empirical support to the previously debated topic: \textit{does pre-training on code boost reasoning capabilities?} Our results affirm this, at least within the domain of mathematical reasoning.

    \item Contrary to prevailing practice, incorporating arXiv publications in training—a common strategy among models with mathematical tasks—does not yield measurable gains across the suite of mathematical benchmarks evaluated herein.
\end{itemize}

\textbf{Exploration and Analysis of Reinforcement Learning}

\begin{itemize}[topsep=0pt]
    \item We present Group Relative Policy Optimization (GRPO), a resource-efficient and robust reinforcement learning approach. By bypassing the critic model and instead deriving the baseline directly from group scores, GRPO substantially reduces the computational requirements compared to Proximal Policy Optimization (PPO).
    \item Our results show that employing GRPO yields notable improvements in our instruction-tuned model DeepSeekMath-Instruct, relying exclusively on instruction-tuning data. Additionally, the reinforcement learning phase yields consistent gains in handling tasks outside of the training domain.
    \item We articulate a cohesive framework that encompasses various strategies including RFT, DPO, PPO, and GRPO. Through a broad suite of experiments—such as contrasting online with offline learning, comparing outcome supervision to process supervision, and evaluating both single-turn and iterative reinforcement learning—we systematically analyze the key factors underlying this unified view.
    \item Leveraging this comprehensive paradigm, we investigate the mechanisms driving the effectiveness of reinforcement learning and propose several avenues for further improving reinforcement learning for LLMs.
\end{itemize}

\subsection{Summary of Evaluations and Metrics}

\begin{itemize}[topsep=0pt]
    \item \textbf{English and Chinese Mathematical Reasoning}: We thoroughly evaluate our models using both English and Chinese mathematical benchmarks spanning primary through tertiary education levels. For English, we utilize datasets such as GSM8K \citep{gsm8k}, MATH \citep{MATH}, SAT \citep{llemma}, OCW Courses \citep{minerva}, and MMLU-STEM \citep{mmlu}. Chinese evaluations include MGSM-zh \citep{mgsm}, CMATH \citep{wei2023cmath}, Gaokao-MathCloze \citep{agieval}, and Gaokao-MathQA \citep{agieval}. We assess model performance in producing comprehensive written solutions unaided by external tools, as well as their capacity to solve mathematical tasks with Python.

    DeepSeekMath-Base demonstrates performance on par with the proprietary Minerva 540B \citep{minerva} on English tasks, and outperforms all open-source foundational models, including Mistral 7B \citep{mistral} and Llemma-34B \citep{llemma}, often by a notable margin regardless of prior mathematical pre-training. Its strong showing on Chinese benchmarks appears to result from incorporating high-quality multilingual pre-training data rather than restricting to English-only sources as in earlier work \citep{minerva,llemma}. Instruction-tuning and reinforcement learning lead to the advanced models DeepSeekMath-Instruct and DeepSeekMath-RL, which achieve outstanding accuracy—surpassing 50\% on the highly competitive MATH dataset—a first within the open-source community.
\end{itemize}

\item \textbf{Formal Mathematics}: DeepSeekMath-Base is assessed on the informal-to-formal theorem proving task introduced in \citep{dsp_proof}, utilizing the miniF2F benchmark \citep{minif2f} and employing Isabelle \citep{isabelle} as the designated proof assistant. The model exhibits impressive few-shot autoformalization capabilities.

\item \textbf{Natural Language Understanding, Reasoning, and Code}: To thoroughly evaluate the general linguistic comprehension, reasoning abilities, and coding proficiency of DeepSeekMath-Base, we benchmark it on the Massive Multitask Language Understanding (MMLU) assessment \citep{mmlu}, which includes 57 multiple-choice questions across a wide range of fields; BIG-Bench Hard (BBH) \citep{bbh}, featuring 23 tasks that largely demand complex, multi-step reasoning; and both HumanEval \citep{codex} and MBPP \citep{mbpp}, which are standard benchmarks for code-oriented language models. Pre-training with mathematical content enhances the model’s performances in both language understanding and reasoning.

\section{Math Pre-Training}

\subsection{Data Collection and Decontamination} This section details our methodology for developing the DeepSeekMath Corpus based on data from Common Crawl. As illustrated in Figure \ref{fig:data_collect}, we establish an iterative pipeline for systematically extracting a large-scale mathematical dataset from Common Crawl, beginning with a high-quality but limited seed dataset of math content. Notably, the procedure described here is adaptable for other specialized domains, such as programming.

Our process starts with OpenWebMath \citep{openwebmath}, which offers a reliable compilation of mathematical web texts, chosen as the initial seed set. Using OpenWebMath, we build a fastText classifier \citep{joulin2016fasttext} to retrieve additional web pages that are similar in content and style to OpenWebMath. The classifier is trained on a randomly sampled set of 500,000 positive instances from the seed and 500,000 negative instances from Common Crawl. We utilize the public fastText implementation\footnote{\url{https://fasttext.cc}}, configuring it with a 256-dimensional embedding, a learning rate of 0.1, word n-gram maximum length of 3, minimum word count threshold of 3, and 3 training epochs. To efficiently condense the original Common Crawl data, we implement both exact and near-duplicate URL removal, reducing the collection to 40 billion HTML pages. The trained fastText model is then employed to select math-related pages from this deduplicated dataset. To exclude inferior quality content, the retrieved documents are ranked by their fastText-predicted scores, and only the top-scoring subset is retained. The chosen quantity of data is empirically evaluated via pre-training trials on datasets consisting of the top 40B, 80B, 120B, and 160B tokens, with the top 40B tokens being selected during the first round.

Following the initial round of data acquisition, a substantial portion of mathematical web pages remain uncollected, primarily due to the limited variety present in the positive training examples used for the fastText model. To address this issue and enhance model performance, we identify further mathematical web sources to supplement the initial seed dataset. We begin by segmenting the complete Common Crawl dataset into mutually exclusive domains, with each domain defined as all web pages sharing an identical base URL. For each domain, we compute the proportion of pages retrieved during the first collection phase. Domains where more than 10\% of pages were collected are considered likely to contain mathematical material (such as \url{mathoverflow.net}). Within these domains, we proceed to manually label URLs that are known to host mathematical content (for instance, \url{mathoverflow.net/questions}). Any uncollected web pages linked to these URLs are incorporated into the seed set. This method allows us to acquire a greater diversity of positive samples, thereby equipping the fastText model to recover a broader range of mathematical content in subsequent cycles. After performing four such collection rounds, we gather a total of 35.5M mathematical web pages, encompassing 120B tokens. By the fourth cycle, we observe that approximately 98\% of the data had already been amassed during the third, which leads us to terminate further data acquisition.

To prevent contamination of evaluation benchmarks, we adopt the filtering strategy outlined by \cite{deepseek-coder}, excluding from our math training corpus all web pages that include content overlapping with English mathematical benchmarks like GSM8K \citep{gsm8k} and MATH \citep{MATH}, as well as Chinese datasets such as CMATH \citep{wei2023cmath} and AGIEval \citep{agieval}. The specific criterion applied is as follows: any segment of text in our dataset is removed if it contains an exact match to any 10-gram substring found in the benchmarks. In cases where benchmark content is shorter than 10 grams but at least 3 grams long, we use an exact match protocol to excise any affected web pages.

\subsection{Validating the Quality of the DeepSeekMath~Corpus}

We carry out pre-training experiments to evaluate how DeepSeekMath~Corpus compares to recent mathematics-focused corpora:

\begin{itemize}[topsep=0pt]
    \item \textbf{MathPile} \citep{mathpile}: This is a collection of 8.9B tokens curated from sources such as textbooks, Wikipedia, ProofWiki, CommonCrawl, StackExchange, and arXiv, with arXiv comprising more than 85\% of the data;
    \item \textbf{OpenWebMath} \citep{openwebmath}: Derived from CommonCrawl and filtered for mathematical relevance, totaling 13.6B tokens;
    \item \textbf{Proof-Pile-2} \citep{llemma}: Combines OpenWebMath, AlgebraicStack (10.3B tokens of mathematical code), and arXiv manuscripts (28.0B tokens). For our experiments, we apply the arXiv:Web:Code ratio of 2:4:1, consistent with \cite{llemma}.
\end{itemize}

\subsubsection{Training Setting}
\label{sec:quality-policy}
We perform specialized mathematical training on a general-purpose language model consisting of 1.3B parameters, architecturally aligned with the DeepSeek LLMs \citep{deepseek-llm}, referred to as DeepSeek-LLM 1.3B. Each distinct mathematical dataset is used to train a separate model instance over 150B tokens. The entirety of the training experiments utilize the HAI-LLM \citep{haillm} framework, noted for its efficiency and low resource demands. Adhering to the methodology employed in DeepSeek LLMs, optimization is conducted via AdamW \citep{adamW}, setting $\beta_1=0.9$, $\beta_2=0.95$, and $\mathrm{weight\_decay}=0.1$. The learning rate follows a staged schedule: it linearly ramps up to its peak after 2,000 warmup iterations, then falls to 31.6\% of the maximum by the 80\% mark of training completion, and further drops to 10.0\% by 90\% progress. The maximum learning rate is established at 5.3e-4. Each training step processes a batch of 4M tokens, with a context window of 4K tokens.

\subsubsection{Evaluation Results}

\textbf{The DeepSeekMath~Corpus is distinguished by its high quality, substantial scale, and comprehensive multilingual scope.}

\begin{itemize}[topsep=0pt]
    \item \textbf{High-quality}:
    We assess downstream task effectiveness across eight mathematical evaluation sets using few-shot chain-of-thought prompting \cite{cot}. Table \ref{tab:corpora_comparison} reveals a distinct advantage for models trained on the DeepSeekMath~Corpus. Furthermore, as shown in Figure \ref{fig:corpora_comparisons}, models trained on the DeepSeekMath Corpus consistently outperform those trained on Proof-Pile-2 at the 50B token milestone (corresponding to one epoch on Proof-Pile-2), signifying a higher average data quality within the DeepSeekMath Corpus.
    \item \textbf{Multilingual}:
    The DeepSeekMath~Corpus contains mathematical data across several languages, with English and Chinese as the most prevalent. Analysis in Table \ref{tab:corpora_comparison} shows that exposure to DeepSeekMath~Corpus data bolsters mathematical reasoning in both English and Chinese, whereas most available mathematical datasets—primarily in English—offer limited benefit and may even impair model capability on Chinese mathematical tasks.
    \item \textbf{Large-scale}:
    DeepSeekMath~Corpus is vastly larger than pre-existing mathematical corpora. As visualized in Figure \ref{fig:corpora_comparisons}, training DeepSeek-LLM 1.3B with the DeepSeekMath~Corpus results in a more pronounced learning curve and sustained performance gains. In contrast, smaller baseline corpora, already extensively cycled during training, quickly drive model results to an early plateau.
\end{itemize}

\subsection{Training and Evaluating DeepSeekMath-Base 7B}

This section details the development of DeepSeekMath-Base 7B, a foundational model characterized by its advanced reasoning capabilities in mathematics. The model builds on DeepSeek-Coder-Base-v1.5 7B \citep{deepseek-coder}, from which it inherits its initial parameters, and undergoes further training over 500B tokens. The training data composition is as follows: 56\% originates from the DeepSeekMath Corpus, 4\% is sourced from AlgebraicStack, 10\% from arXiv papers, 20\% from code repositories on Github, while the remaining 10\% comprises natural language samples from Common Crawl in both English and Chinese. The majority of training follows the protocol outlined in Section \ref{sec:quality-policy}, with two exceptions: the peak learning rate is increased to 4.2e-4 and the batch size is adjusted to 10M tokens.

We systematically evaluate the mathematical proficiency of DeepSeekMath-Base 7B, emphasizing its capacity to generate comprehensive mathematical solutions independently, its competence in solving mathematical tasks utilizing tools, and its performance in formal theorem verification. Additionally, a broader evaluation is conducted, which explores the model’s abilities in natural language processing, logical reasoning, and computer programming.

\paragraph{Mathematical Problem Solving with Step-by-Step Reasoning}

To assess DeepSeekMath-Base’s mathematical problem-solving prowess, we utilize few-shot chain-of-thought prompting \citep{cot} across eight benchmarks in both English and Chinese. These evaluation benchmarks include quantitative reasoning datasets—such as GSM8K \citep{gsm8k}, MATH \citep{MATH}, and CMATH \citep{wei2023cmath}—as well as multiple-choice assessments like MMLU-STEM \citep{mmlu} and Gaokao-MathQA \citep{agieval}. Together, they span a broad spectrum of mathematical disciplines and difficulties, ranging from elementary to undergraduate levels.

As indicated in Table \ref{tab:base_math_cot}, DeepSeekMath-Base 7B attains the highest scores across all eight evaluated benchmarks when compared to other open-source base models, including the popular general-purpose model Mistral 7B \citep{mistral} and the more recently introduced Llemma 34B \citep{llemma}, which was specifically trained on Proof-Pile-2 \citep{llemma} for mathematics. Particularly noteworthy is DeepSeekMath-Base's performance on the competition-level MATH benchmark, where it exceeds all open-source base models by a margin exceeding 10 percentage points, and even surpasses Minerva 540B \citep{minerva}, a proprietary model of substantially greater scale—77 times larger—that extends PaLM \citep{palm} with additional training on mathematical content.

\paragraph{Mathematical Problem Solving with Tool Use} 

For program-based mathematical reasoning, we conduct assessments on GSM8K and MATH datasets, employing few-shot program-of-thought prompting techniques \citep{pot,pal}. The models are provided with prompts that require the generation of Python code to resolve each mathematical problem, with access to computational libraries such as \textit{math} and \textit{sympy} to manage complex calculations. The correctness of the models is measured by evaluating the output produced after executing the generated code. Results presented in Table \ref{tab:base_math_pot_proof} reveal that DeepSeekMath-Base 7B surpasses the previous state-of-the-art model Llemma 34B in these settings.

\paragraph{Formal Mathematics}

Automating formal proof construction enhances both the correctness and efficiency of mathematical reasoning, an area that has seen increased scholarly attention. DeepSeekMath-Base 7B is tested on the informal-to-formal proof generation task outlined in \citep{dsp_proof}, wherein the goal is to synthesize a formal proof based on an informal statement, its formal version, and a provided informal proof. The evaluation utilizes miniF2F \citep{minif2f}, a challenging Olympiad-style mathematics benchmark, where few-shot prompting guides the model to generate formal proofs in Isabelle for each problem. Following the procedure in \cite{dsp_proof}, model outputs are proof sketches, subsequently completed using Sledgehammer \citep{sledgehammer}, an established automated prover. Table \ref{tab:base_math_pot_proof} demonstrates that DeepSeekMath-Base 7B is highly effective at automatic proof formalization.

\paragraph{Natural Language Understanding, Reasoning, and Code}

To assess abilities beyond mathematics, we evaluate the model’s competence in natural language understanding on MMLU \citep{mmlu}, logical reasoning on BBH \citep{bbh}, and programming skills using HumanEval \citep{codex} and MBPP \citep{mbpp}. Table \ref{tab:understanding_reasoning_code} shows that DeepSeekMath-Base 7B makes substantial gains over its predecessor, DeepSeek-Coder-Base-v1.5 \citep{deepseek-coder}, in both MMLU and BBH benchmarks, which highlights the value of mathematics-centered training for overall language comprehension and logical reasoning. Furthermore, thanks to continued inclusion of code data in training, DeepSeekMath-Base 7B preserves the code generation performance observed in DeepSeek-Coder-Base-v1.5 for both code evaluation tasks. Taken together, DeepSeekMath-Base 7B consistently outperforms the general-purpose Mistral 7B \citep{mistral} across reasoning and programming benchmarks.

\section{Supervised Fine-Tuning}

\subsection{SFT Data Curation}

We develop an instruction-tuning dataset for mathematics that encompasses both English and Chinese problem statements, drawing from a broad spectrum of mathematical topics and spanning multiple difficulty levels. Each problem is paired with a detailed solution, utilizing formats such as chain-of-thought (CoT) \citep{cot}, program-of-thought (PoT) \citep{pot,pal}, and tool-integrated reasoning \citep{tora}. The final dataset comprises 776K training instances.

\begin{itemize}[topsep=0pt]
    \item \textbf{English mathematical datasets}: Our data includes GSM8K and MATH, enhanced with tool-integrated answer annotations. Additionally, we incorporate selected problems from MathInstruct \citep{MathInstruct}, and the Lila-OOD training corpus \citep{lila} where solutions are expressed using CoT or PoT methods. The English subset represents a variety of mathematical disciplines, such as algebra, probability, number theory, calculus, and geometry.
    \item \textbf{Chinese mathematical datasets}: For Chinese, we curate K-12 mathematics problems that address 76 different subfields, such as linear equations, providing corresponding solutions annotated in both CoT and tool-integrated styles.
\end{itemize}

\subsection{Training and Evaluating DeepSeekMath-Instruct 7B}

Here, we present DeepSeekMath-Instruct 7B, which is further fine-tuned for mathematical instructions based on DeepSeekMath-Base. To utilize long sequences, training samples are concatenated at random until the context window reaches up to 4K tokens. Training is conducted over 500 iterations, with a batch size of 256 and a fixed learning rate set to 5e-5.

We assess the mathematical reasoning abilities of our models, both in pure reasoning mode and when external tools are available, across four quantitative reasoning benchmarks in both English and Chinese. Our evaluation includes comparisons with top-performing contemporary models:

\begin{itemize}[topsep=0pt]
    \item \textbf{Closed-source models}: This category features (1) the GPT series, notably GPT-4 \citep{gpt4} and the GPT-4 Code Interpreter~\footnote{\url{https://openai.com/blog/chatgpt-plugins\#\#code-interpreter}}, (2) Gemini Ultra and Pro \citep{gemini}, (3) Inflection-2 \citep{inflection-2}, (4) Grok-1~\footnote{\url{https://x.ai/model-card}}, alongside models newly made available by Chinese technology firms such as (5) Baichuan-3~\footnote{\url{https://www.baichuan-ai.com}}, and (6) the latest release from the GLM series, GLM-4~\footnote{\url{https://open.bigmodel.cn/dev/api\#glm-4}} \citep{glm}.
\end{itemize}

Most of these models are designed for broad applications and have typically been subjected to various alignment processes.

\item \textbf{Open-source models} consist of: (1) DeepSeek-LLM-Chat 67B \citep{deepseek-llm}, (2) Qwen 72B \citep{qwen}, (3) SeaLLM-v2 7B \citep{seallm}, and (4) ChatGLM3 6B \citep{chatglm3}, which are all general-purpose; in addition to several models with targeted mathematical capabilities such as (5) InternLM2-Math 20B~\footnote{\url{https://github.com/InternLM/InternLM-Math}}, which extends InternLM2 through dedicated mathematics pre-training and subsequent instruction tuning, (6) Math-Shepherd-Mistral 7B, applying PPO training \citep{schulman2017proximal} to Mistral 7B \citep{mistral} by leveraging a process-supervised reward model, (7) the WizardMath suite \citep{wizardmath} which enhances the mathematical reasoning abilities of Mistral 7B and Llama-2 70B \citep{llama2} via evolve-instruct—a variation of instruction tuning utilizing AI-evolved tasks—and PPO training, with major training datasets drawn from GSM8K and MATH, (8) MetaMath 70B \citep{metamath}, where Llama-2 70B is further fine-tuned using an expanded version of GSM8K and MATH, (9) ToRA 34B \cite{tora}, adapting CodeLlama 34B for mathematical reasoning with integrated tool usage, and (10) MAmmoTH 70B \citep{MathInstruct}, based on Llama-2 70B that has undergone instruction tuning using MathInstruct.

As detailed in Table \ref{tab:sft_rl_math}, when tool access is not permitted during evaluation, DeepSeekMath-Instruct 7B delivers robust stepwise reasoning capabilities. Remarkably, on the challenging MATH dataset, our model outperforms all available open-source alternatives as well as most closed-source models (such as Inflection-2 and Gemini Pro) by a minimum margin of 9 percentage points. This superior performance holds even in comparison to models with significantly larger scale (such as Qwen 72B) or those that have benefited from reinforcement learning oriented toward mathematical reasoning (such as WizardMath-v1.1 7B). While DeepSeekMath-Instruct attains results comparable to proprietary Chinese models like GLM-4 and Baichuan-3, it does not yet reach the performance of GPT-4 and Gemini Ultra on MATH.

In the evaluation context where models are permitted to combine natural language reasoning and the use of external programs as problem-solving tools, DeepSeekMath-Instruct 7B achieves nearly 60\% accuracy on MATH, outpacing all prior open-source competitors. On other evaluation sets, our model performs on par with DeepSeek-LLM-Chat 67B, which previously held state-of-the-art status among open-source models but is an order of magnitude larger.

\section{Reinforcement Learning}

\subsection{Group Relative Policy Optimization}

It has been shown that reinforcement learning (RL) techniques can substantially boost the mathematical reasoning skills of large language models (LLMs) following Supervised Fine-Tuning (SFT) \citep{wang2023math,wizardmath}. This section presents Group Relative Policy Optimization (GRPO), our proposed RL method that achieves both effectiveness and efficiency.

\subsubsection{From PPO to GRPO}

Proximal Policy Optimization (PPO) \citep{schulman2017proximal} is a widely adopted actor-critic RL method utilized in the RL-based fine-tuning phase for LLMs \citep{ouyang2022training}. The PPO algorithm aims to refine the policy by maximizing the following surrogate objective:

\begin{equation}
\footnotesize
    \mathcal{J}_{PPO}(\theta) = \mathbb{E}{[q \sim P(Q), o \sim \pi_{\theta_{old}}(O|q)]} \frac{1}{|o|} \sum_{t=1}^{|o|} \min \left[ \frac{\pi_\theta(o_{t} | q, o_{<t})}{\pi_{\theta_{old}}(o_{t} | q, o_{<t})} A_{t}, \text{clip} \left( \frac{\pi_\theta(o_{t} | q, o_{<t})}{\pi_{\theta_{old}}(o

In this context, $\pi_{\theta}$ and $\pi_{\theta_{old}}$ denote the present and prior policy models, respectively, while $q$ and $o$ represent questions and responses sampled from the dataset and from $\pi_{\theta_{old}}$. The parameter $\varepsilon$ is a hyperparameter associated with the clipping mechanism in PPO, designed to enhance training stability. The advantage term $A_t$ is calculated through Generalized Advantage Estimation (GAE) \citep{gae}, leveraging the sequence of rewards $\{r_{\ge t}\}$ and a value function $V_{\psi}$ learned in parallel. As a consequence, PPO requires concurrent training of a value function. To counteract potential over-optimization of the reward model, it is standard practice to incorporate a per-token KL divergence penalty from a reference model into the reward at every token position \citep{ouyang2022training}, as follows:

\begin{equation}
    r_{t} = r_\varphi(q, o_{\le t}) - \beta \log\frac{\pi_{\theta}(o_{t}|q, o_{<t})}{\pi_{ref}(o_{t}|q, o_{<t})},
\label{eq:PPO-reward}
\end{equation}

Here, $r_\varphi$ denotes the reward model, $\pi_{ref}$ is the reference policy (commonly the original SFT model), and $\beta$ serves as the KL penalty weighting factor.

Utilizing a value function in PPO often means maintaining a model roughly the same scale as the policy, which significantly increases the memory and compute requirements. Furthermore, in RL training, the value function's role is to serve as a baseline for variance reduction when estimating advantages. However, since in large language models typically only the final token receives a reward label from the reward model, reliably training a per-token value function becomes challenging. To mitigate these issues, we introduce Group Relative Policy Optimization (GRPO), illustrated in Figure \ref{fig:grpo}, which dispenses with explicit value function estimation as required in PPO. Instead, GRPO computes the baseline using the average reward from multiple responses sampled for each question. Specifically, for any question $q$, GRPO generates a batch of outputs $\{o_1, o_2, \cdots, o_G\}$ using $\pi_{\theta_{old}}$ and updates the policy by maximizing:

\begin{equation}
\footnotesize
\begin{split}
    \mathcal{J}_{GRPO}(\theta) &= \mathbb{E}{[q \sim P(Q), \{o_i\}_{i=1}^G \sim \pi_{\theta_{old}}(O|q)]}  \\
    & \frac{1}{G}\sum_{i=1}^G\frac{1}{|o_i|} \sum_{t=1}^{|o_i|} \left\{ \min \left[ \frac{\pi_\theta(o_{i,t} | q, o_{i,<t})}{\pi_{\theta_{old}}(o_{i,t} | q, o_{i,<t})} \hat{A}_{i,t}, \text{clip} \left( \frac{\pi_\theta(o_{i,t} | q, o_{i,<t})}{\pi_{\theta_{old}}(o_{i,t} | q, o_{i,<t})}, 1 - \varepsilon, 1 + \varepsilon \right)  \hat{A}_{i,t} \right] - \beta \mathbb{D}_{KL}\left[\pi_{\theta} || \pi_{ref}\right]\right\} ,
\end{split}
\label{eq:GRPO-obj}
\end{equation}

In this formulation, $\varepsilon$ and $\beta$ are hyperparameters, and $\hat{A}_{i,t}$ represents the advantage, computed based solely on the relative rewards within each group—further details are provided in the following sections. By calculating advantages within a group of outputs for the same question, GRPO aligns with the inherently comparative design of reward models, which are generally trained using paired comparisons on shared prompts. Importantly, GRPO applies KL regularization by directly including the divergence term between the updated policy and the reference model in the loss, bypassing the need to adjust the reward definition (as in (\ref{eq:PPO-reward})) and simplifying computation of $\hat{A}_{i,t}$. Unlike the KL penalty in (\ref{eq:PPO-reward}), we estimate the KL divergence through an unbiased estimator from \citep{kl_approx}:

\begin{equation}
\small
    \mathbb{

\begin{algorithm}[t]
  \small
  \caption{Iterative Group Relative Policy Optimization}
  \textbf{Input} initial policy model $\pi_{\theta_{\text{init}}}$; reward models $r_\varphi$; task prompts $\mathcal{D}$; 
  hyperparameters $\varepsilon$, $\beta$, $\mu$
  \begin{algorithmic}[1]
    \State Initialize policy model: $\pi_\theta \leftarrow \pi_{\theta_{\text{init}}}$
    \For{iteration = 1, \dots, I}
       \State Assign $\pi_{ref} \leftarrow \pi_{\theta}$ as the reference model
      \For{step = 1, \dots, M}
      \State Draw a mini-batch $\mathcal{D}_b$ from $\mathcal{D}$
      \State Set $\pi_{\theta_{old}} \leftarrow \pi_{\theta}$ to preserve current policy parameters 
      \State For every question $q \in \mathcal{D}_b$, generate $G$ completions $\{o_i\}_{i=1}^G$ via $\pi_{\theta_{old}} (\cdot \mid q)$ 
      \State Obtain rewards $\{r_i\}_{i=1}^{G}$ for each $o_i$ by evaluating with $r_{\varphi}$ 
      \State Employ group relative advantage estimation to calculate $\hat{A}_{i,t}$ for the $t$-th token in each $o_i$
      \For{GRPO iteration = 1, \dots, $\mu$}
        \State Refine policy model $\pi_{\theta}$ by maximizing the GRPO objective as described in Equation \ref{eq:GC-GRPO}
      \EndFor
    \EndFor 
    \State Continuously improve $r_\varphi$ by applying a replay training strategy. 
    \EndFor 
  \end{algorithmic}
  \textbf{Output} $\pi_\theta$
  \label{alg:iter-grpo}
\end{algorithm}

\subsubsection{Outcome Supervision RL with GRPO}
In this approach, for any given question $q$, $G$ output samples $\{o_1, o_2, \cdots, o_G\}$ are generated using the previous policy model $\pi_{\theta_{old}}$. Each output is evaluated by a reward model, resulting in a set of rewards $\mathbf{r} = \{r_1, r_2, \cdots, r_G\}$. These rewards are standardized by removing the group mean and dividing by their standard deviation. Under outcome supervision, the normalized reward is only provided at the conclusion of each response $o_i$, and every token in $o_i$ is assigned an advantage $\hat{A}_{i, t}$ equal to this normalized reward: $\hat{A}_{i, t} = \widetilde{r}_i = \frac{r_i- {\rm mean}(\mathbf{r})}{{\rm std}(\mathbf{r})}$. The policy is then improved by maximizing the objective in Equation~(\ref{eq:GRPO-obj}).

\subsubsection{Process Supervision RL with GRPO}
While outcome supervision grants a terminal reward for each output, this may lack granularity and be suboptimal in supervising the policy on sophisticated mathematical problems. Building on the methods of \cite{wang2023math}, we also investigate process supervision, which delivers rewards at the completion of each individual reasoning step. Formally, provided question $q$ and sampled sequences $\{o_1, o_2, \cdots, o_G\}$, a process reward model scores every step, producing a set of stepwise rewards: $\mathbf{R} = \{ \{r_1^{index(1)},\cdots,r_1^{index(K_1)}\}, \cdots,  \{r_G^{index(1)},\cdots,r_G^{index(K_G)}\} \}$. Here, $index(j)$ locates the final token of the $j$-th step, with $K_i$ as the count of steps in output $i$. As before, these stepwise rewards are normalized across the group: $\widetilde{r}_i^{index(j)} = \frac{r_i^{index(j)} - {\rm mean(\mathbf{R})}}{{\rm std(\mathbf{R})}}$. 
For process supervision, the advantage assigned to each token equals the total normalized rewards from all subsequent reasoning steps, that is, $\hat{A}_{i, t} = \sum_{index(j) \ge t} \widetilde{

\subsubsection{Iterative RL with GRPO} Throughout the reinforcement learning procedure, the previously trained reward model may no longer provide adequate supervision for the evolving policy model. To address this, we examine an iterative approach to RL using GRPO. As illustrated in Algorithm \ref{alg:iter-grpo}, our iterative GRPO strategy involves generating updated reward model training datasets derived from samples produced by the current policy model. The reward model is continually refined through a replay mechanism, integrating 10\% of past data into each training cycle. Subsequently, we use the latest policy model as the reference model, and proceed to train the policy model anew using the updated reward model.

\subsection{Training and Evaluating DeepSeekMath-RL}

We perform reinforcement learning based on the DeepSeekMath-Instruct 7B architecture. The RL training corpus comprises approximately 144K chain-of-thought questions related to GSM8K and MATH, sourced from the SFT dataset. Other SFT questions are omitted to specifically assess the effect of reinforcement learning on benchmarks not represented during the RL stage. Construction of the reward model’s training dataset follows the methodology outlined in \citep{wang2023math}. The initial reward model is trained on DeepSeekMath-Base 7B with a learning rate of 2e-5. For training the policy model via GRPO, we adopt a learning rate of 1e-6, and the KL penalty coefficient is fixed at 0.04. Each question is associated with $64$ sampled completions. The sequence length is capped at 1024 tokens, and each training batch also contains 1024 samples. Following every exploration iteration, the policy model is updated only once. For evaluation, DeepSeekMath-RL 7B is benchmarked in the same fashion as DeepSeekMath-Instruct 7B. Tasks involving GSM8K and MATH with chain-of-thought reasoning are treated as in-domain, while all remaining benchmarks serve as out-of-domain evaluations.

Table \ref{tab:sft_rl_math} presents the comparative results of both open-source and closed-source models, evaluated on English and Chinese benchmarks employing chain-of-thought and tool-assisted reasoning approaches. The key findings include: (1) DeepSeekMath-RL 7B achieves accuracy rates of 88.2\% on GSM8K and 51.7\% on MATH by leveraging chain-of-thought reasoning. These results not only exceed those of all other open-source models ranging from 7B to 70B parameters but also outperform the majority of closed-source counterparts. (2) Notably, DeepSeekMath-RL 7B’s training regime is limited to chain-of-thought instruction tuning using only the GSM8K and MATH datasets, building directly upon DeepSeekMath-Instruct 7B. Despite this limited data exposure, it consistently surpasses DeepSeekMath-Instruct 7B on all assessment metrics, highlighting the benefits introduced by reinforcement learning.

\section{Discussion}
This section reports our insights gained through pre-training and reinforcement learning experiments.

\subsection{Insights from Pre-Training}

We begin by detailing the insights obtained during pre-training. Unless noted otherwise, our discussion adheres to the experimental configuration described in Section \ref{sec:quality-policy}. It is also important to clarify that references to the DeepSeekMath Corpus in this context pertain to an 89B-token subset assembled in the second round of dataset compilation.

\subsubsection{Code Training Enhances Mathematical Reasoning}

A widely held—though not yet rigorously verified—notion in the community is that exposure to code improves a model's reasoning capabilities. Here, we provide partial empirical support for this claim within the context of mathematical reasoning: code training augments models' mathematical reasoning skills, regardless of whether tool integration is employed.

To investigate the impact of code training on mathematical reasoning, we examined and compared outcomes across a two-stage and a one-stage training pipeline:

\textbf{Two-Stage Training}

\begin{itemize}[topsep=0pt]
    \item \textbf{Code Training for 400B Tokens $\rightarrow$ Math Training for 150B Tokens}: DeepSeek-LLM 1.3B is first pretrained on 400B code tokens, followed by an additional phase of 150B math tokens.
    \item \textbf{General Training for 400B Tokens $\rightarrow$ Math Training for 150B Tokens}: To serve as a baseline, we replace the initial code token training with 400B general tokens (sourced from DeepSeek-AI’s extensive general text dataset), then continue with the same 150B math token phase. This enables a direct comparison regarding the impact of pretraining on code tokens versus general tokens for enhancing mathematical reasoning capabilities.
\end{itemize}

\textbf{One-Stage Training}

\begin{itemize}[topsep=0pt]
    \item \textbf{Math Training for 150B Tokens}: DeepSeek-LLM 1.3B is trained from the outset solely on 150B math tokens.
    \item \textbf{Training on a mixture of 400B Code Tokens and 150B Math Tokens}: When math training is conducted after a dedicated code training phase, there is a notable decline in code task performance. Therefore, we evaluate whether simultaneously combining 400B code tokens with 150B math tokens in a single-stage training can reinforce mathematical reasoning while reducing the extent of catastrophic forgetting.
\end{itemize}

\paragraph{Results}
Table \ref{tab:code-math} and Table \ref{tab:code-math-general-eval} provide empirical evaluations under various training protocols.

Exposure to code tokens during training significantly supports program-aided mathematical reasoning, as observed in both staged and joint training paradigms. Table \ref{tab:code-math} indicates that pretraining exclusively on code tokens in the two-stage setup confers notable proficiency in resolving GSM8K and MATH problems with Python. Augmenting with math token finetuning in the subsequent stage provides further gains. Notably, in the single-stage approach where code and math tokens are interleaved, the adverse effects of catastrophic forgetting apparent in sequential training are largely prevented. This strategy enhances both coding task performance (as detailed in Table \ref{tab:code-math-general-eval}) and supports strong program-based mathematical problem-solving (see Table \ref{tab:code-math}).

Furthermore, pretraining on code is also advantageous for math reasoning tasks that do not involve external tools. In the two-stage protocol, early code exposure yields moderate improvements and positively influences the efficacy of later math token training, ultimately achieving superior outcomes. Conversely, a single-stage blend of code and math tokens appears to inhibit performance on pure mathematical reasoning. This could be attributed to the restricted representational capacity of DeepSeek-LLM 1.3B, making it challenging for the model to assimilate both code and math content in a single unified training phase.

\subsubsection{ArXiv Papers Appear Unhelpful for Advancing Mathematical Reasoning}

ArXiv papers have traditionally been included as a significant part of mathematical pre-training datasets \citep{minerva,gpt-f,llemma,mathpile}. Yet, comprehensive studies scrutinizing their specific contribution to mathematical reasoning are lacking. Somewhat unexpectedly, our experimental findings suggest that arXiv papers offer limited, if any, benefit in enhancing mathematical reasoning abilities. To investigate this, we utilized a range of model scales—specifically DeepSeek-LLM 1.3B and DeepSeek-Coder-Base-v1.5 7B \citep{deepseek-coder}—paired with distinct versions of arXiv corpora subjected to different preprocessing routines:

\begin{itemize}[topsep=0pt]
    \item \textbf{MathPile} \citep{mathpile}: An 8.9-billion-token collection extensively cleaned and filtered via heuristic methods, with more than 85\% content sourced from scientific arXiv publications;
    \item \textbf{ArXiv-RedPajama} \citep{redpajama}: The full set of arXiv LaTeX documents stripped of preambles, comments, macros, and references, comprising 28.0 billion tokens.
\end{itemize}

Within our study, we trained DeepSeek-LLM 1.3B over 150 billion tokens and DeepSeek-Coder-Base-v1.5 7B over 40 billion tokens on each respective arXiv dataset. Our observations indicate that relying solely on arXiv texts yields negligible or even negative effects on mathematical reasoning benchmarks of varying complexity. In particular, both models trained exclusively on these corpora failed to show significant improvements, sometimes exhibiting declines on several standardized evaluations, including quantitative reasoning tasks such as GSM8K and MATH (see Table \ref{tab:arxiv-cot}), multiple-choice assessments like MMLU-STEM (see Table \ref{tab:arxiv-cot}), and formal mathematics suites such as miniF2F (see Table \ref{tab:arxiv_atp}).

Nevertheless, the scope of these conclusions is not absolute, and caution is warranted. Our current research does not address:

\begin{itemize}[topsep=0pt]
    \item How arXiv token data affects specialized mathematics tasks outside our benchmarks, such as transforming formal statements or proofs into their informal counterparts (informalization);
    \item The potential interactions or improvements when arXiv data are merged with other types of training corpora;
    \item Whether positive effects from arXiv papers may only become evident at model sizes larger than those tested here.
\end{itemize}

Accordingly, these open questions suggest that more thorough investigations are necessary, which we plan to pursue in subsequent work.

\subsection{Insights into Reinforcement Learning}
\subsubsection{Toward a Unified Training Framework} In what follows, we propose a comprehensive analytical framework to examine a spectrum of training algorithms, including SFT, RFT, DPO, PPO, and GRPO, and empirically investigate how various components of this framework affect performance. In general, the gradient with respect to the parameters $\theta$ for a given training method may be represented as:

\begin{equation}
    \nabla_{\theta}\mathcal{J}_{\textcolor{red}{\mathcal{A}}}(\theta) = \mathbb{E}[\underbrace{(q,o) \sim \textcolor{red}{\mathcal{D}}}_{Data \ Source}]\left( \frac{1}{|o|} \sum_{t=1}^{|o|}  \underbrace{GC_{{\mathcal{A}}}(q, o, t, \textcolor{red}{\pi_{{rf}}})}_{Gradient \ Coefficient}  \nabla_{\theta}\log \pi_{\theta}(o_t | q, o_{<t})\right).
\label{eq:objective}
\end{equation}

This formulation consists of three principal elements: 1) the \textit{Data Source $\mathcal{D}$}, specifying the dataset employed for model training; 2) the \textit{Reward Function $\pi_{{rf}}$}, serving as the generator of reward signals for learning; and 3) the \textit{Algorithm $\mathcal{A}$}, which combines the inputs from the data and reward signal into a gradient coefficient $GC$, thereby guiding the extent of reward or penalty administered. Utilizing this common framework, we can discuss and compare multiple widely-used methods:

\begin{itemize}
    \item  \textbf{Supervised Fine-tuning (SFT)}: SFT adapts a pretrained model to new tasks by further training it on SFT datasets curated by humans.
    \item \textbf{Rejection Sampling Fine-tuning (RFT)}: RFT involves an additional training stage, where the SFT model is fine-tuned with model outputs sampled from SFT prompts, but only those outputs that have been validated for correctness are retained.
    \item \textbf{Direct Preference Optimization (DPO)}: DPO improves the SFT model by fine-tuning it on expanded sets of model responses, applying a pairwise DPO objective to incorporate preference-based learning.
    \item \textbf{Online Rejection Sampling Fine-tuning (Online RFT)}: Unlike standard RFT, Online RFT starts with the SFT model as the initial policy, then continuously refines the policy using outputs dynamically sampled from the latest policy iteration.
    \item \textbf{PPO/GRPO}: These approaches initialize the policy from the SFT model and iteratively strengthen it using outputs produced and sampled from the evolving policy model itself.
\end{itemize}

A summary of the components inherent to these approaches is provided in Table \ref{tab:data_policy}. For an expanded exposition of the derivation process, consult Appendix \ref{app:analysis-rl}.

\paragraph{Analysis of Data Source} We classify the origin of data into two primary types: online sampling and offline sampling. In online sampling, training data is generated from the exploratory actions of the currently updated policy model during training, whereas offline sampling utilizes training data sampled from the outputs of the initial SFT model. Within this framework, both RFT and DPO employ offline sampling strategies, while Online RFT and GRPO rely on online sampling approaches.

Figure \ref{fig:rl-analysis} illustrates that Online RFT achieves substantially higher performance than RFT across two benchmarks. Initially, both Online RFT and RFT demonstrate similar effectiveness; however, as training progresses, Online RFT exhibits a pronounced improvement, showcasing the efficacy of online training. This outcome aligns with expectations: at the outset, the actor and SFT model are nearly indistinguishable, producing sampled data with only minor variations. As training advances, the actor’s sampled data diverges more markedly, and dynamic, real-time data collection yields greater benefits.

\paragraph{Analysis of Gradient Coefficient} The process of updating model parameters involves translating the reward signal into a gradient coefficient. For our experiments, the reward function is categorized as either `Rule' or `Model'. The ‘Rule’ category determines response quality by assessing the accuracy of answers, while the ‘Model’ category leverages a trained reward model to assign scores to responses, with this reward model’s training data grounded in rule-based judgments. Equations \ref{eq:GC-RFT} and \ref{eq:GC-GRPO} underscore a critical divergence between GRPO and Online RFT: GRPO distinctively modulates its gradient coefficient in accordance with the magnitude of the reward assigned by the reward model. This enables the algorithm to provide varied reinforcement or penalization based on the specific reward values. Conversely, Online RFT does not incorporate this mechanism—it does not penalize incorrect outputs and instead applies identical reinforcement to all correct responses, regardless of the degree of correctness.

As illustrated in Figure \ref{fig:rl-analysis}, GRPO outperforms online RFT, underscoring the effectiveness of modifying positive and negative gradient coefficients. Additionally, GRPO+PS achieves better outcomes than GRPO+OS, which demonstrates the advantages of incorporating more granular, step-dependent gradient coefficients. Moreover, we investigate the effects of iterative RL by performing two iterations in our experiments. Figure \ref{fig:iter-rl} reveals that applying RL iteratively substantially boosts performance, particularly during the initial round.

\subsubsection{Why RL Works?} In this study, reinforcement learning is conducted using a subset of instruction tuning data, yielding considerable gains over the instruction-tuned baseline. To better elucidate the underlying mechanism for this improvement, we assess both Pass@K and Maj@K accuracies of the Instruct and RL models across two benchmark tasks. The results, depicted in Figure \ref{fig:maj-pass}, indicate that while RL markedly increases Maj@K accuracy, it does not enhance Pass@K. These observations imply that reinforcement learning strengthens the robustness of the model’s output distribution—specifically, \textbf{the progress seems driven by elevating the correct answer into the TopK predictions, rather than an intrinsic advancement in core capabilities.} This phenomenon echoes the \textbf{misalignment problem} in reasoning highlighted by \citep{wang2023making} for SFT models, and recent works \citep{yuan2023rrhf,song2023preference,wang2023making} suggest that applying preference alignment strategies can mitigate this issue and improve reasoning results.

\subsubsection{How to Achieve More Effective RL?}
\label{sec:effec_rl}
Our experiments reveal that RL is highly effective for mathematical reasoning tasks. We further introduce a unified perspective that accommodates various key training methodologies, interpreting each as a form of direct or approximate RL. According to Equation \ref{eq:objective}, this framework involves three fundamental components: Data Source, Algorithm, and Reward Function. We propose several avenues for future research within these dimensions.

\paragraph{Data Source} The data source underpins every training approach. In RL contexts, it particularly refers to unlabeled problems paired with responses sampled from the policy model. For this work, we leverage only questions from instruction tuning and utilize basic nucleus sampling for output generation. We speculate this restriction may contribute to our RL framework’s improvement being confined to Maj@K performance. Moving forward, we intend to extend our RL pipeline to out-of-distribution prompts and integrate \textbf{advanced sampling (decoding) strategies} such as tree-search-based methods \citep{yao2023tree}. Furthermore, incorporating \textbf{efficient inference techniques} \citep{xia-etal-2023-speculative,leviathan2023fast,kwon2023efficient,xia2024unlocking}—which govern the exploration ability of policy models—will likely be critical for optimizing RL effectiveness.

\paragraph{Algorithms} Algorithms utilize input data along with the reward signal to compute gradient coefficients, which are then applied to adjust the model parameters. According to Equation \ref{eq:objective}, prevailing approaches essentially \textbf{TRUST} the feedback from the reward function to modify the conditional likelihood of specific tokens, either increasing or decreasing it. Yet, there is no guarantee that the reward signal remains accurate across all scenarios, especially in highly intricate tasks. For instance, the PRM800K datasets \citep{lightman2023let}, despite being meticulously labeled by expert annotators, still exhibit an error rate of about 20\% in their annotations\footnote{\url{https://github.com/openai/prm800k/issues/12\#issuecomment-1728491852}}. Therefore, it is crucial to investigate reinforcement learning algorithms designed to withstand unreliable or noisy reward signals. We propose that alignment strategies employing the \textbf{WEAK-TO-STRONG} paradigm \citep{burns2023weak} hold the potential to significantly reshape future learning algorithm designs.

\paragraph{Reward Function} The reward function serves as the foundational signal driving the training process. In reinforcement learning frameworks, this is most often realized through a neural reward model. We identify three primary avenues for advancing reward model research: 1) \textbf{Improving the reward model's generalization capabilities.} A robust reward model must generalize effectively to novel, out-of-distribution queries and more sophisticated decoded responses; otherwise, reinforcement learning might simply consolidate existing distributions in large language models (LLMs) rather than genuinely enhance their competencies; 2) \textbf{Accurately quantifying the uncertainty in the reward model.} Accounting for uncertainty can facilitate the connection between suboptimal reward models and weak-to-strong learning methods; 3) \textbf{Streamlining the creation of high-caliber process reward models} capable of delivering detailed, stage-specific training feedback for complex reasoning tasks \citep{lightman2023let,wang2023math}.

\section{Conclusion, Limitation, and Future Work} 

This work introduces DeepSeekMath, which surpasses all existing open-source models on the MATH benchmark at competition level and achieves results approaching those of proprietary systems. DeepSeekMath~begins with DeepSeek-Coder-v1.5 7B as its foundation and is further continually trained over 500B tokens, utilizing a substantial subset of 120B math-related tokens drawn from Common Crawl. Through a thorough ablation analysis, we demonstrate that web pages serve as a rich source for high-quality mathematical training data, while arXiv contributes less than anticipated. We propose Group Relative Policy Optimization (GRPO), an alternative to Proximal Policy Optimization (PPO), capable of enhancing mathematical reasoning skills with reduced memory overhead. Our experimental findings confirm GRPO's efficacy, even when DeepSeekMath-Instruct 7B already exhibits strong benchmark scores. We also offer a consolidated perspective on various reinforcement learning methodologies and highlight multiple directions for the advancement of reinforcement learning techniques.

Despite DeepSeekMath's~strong performance in quantitative reasoning assessments, it exhibits comparatively weaker results in domains such as geometry and theorem-proving relative to proprietary models. Specifically, in preliminary evaluations, the model failed to address challenges involving triangles and ellipses, possibly pointing to biases introduced during pre-training and fine-tuning data selection. Furthermore, due to limitations in model size, DeepSeekMath~lags behind GPT-4 in few-shot learning: whereas GPT-4 shows marked gains when provided with few-shot prompts, DeepSeekMath~demonstrates little difference between zero-shot and few-shot scenarios. Moving forward, we plan to refine our engineered data selection processes to assemble an even higher-quality pre-training dataset. Moreover, we intend to investigate the research avenues outlined in Section \ref{sec:effec_rl} to further optimize reinforcement learning strategies for large language models.

\end{document}